\section{Introduction}
Agent Skills package reusable instructions with scripts and supporting files.
They can help an agent complete a task consistently, but they also place
third-party content close to tools that can read, change, delete, or disclose
data. Studies have documented malicious Skills and attacks placed directly in
Skill instructions~\cite{liu,skillinject}. Attacks can also depend on the
interaction between apparently reasonable instructions and helper
code~\cite{skillmutator}. Checking a description or searching for a dangerous
tool name is therefore insufficient to establish what a package will do.

The relevant security question depends on the user's task. Sending a report may
be required when the user asks to share it, but excessive when the request is
only to prepare a local copy. Deleting a temporary input may likewise be either
requested cleanup or an unwanted loss of data. A task-aware analysis must link
concrete actions to the requested operation, resource, destination, and side
effects. SkillScope makes this action-in-task view central to Skill
analysis~\cite{skillscope}. It also distinguishes authorization from necessity:
an action can be outside the user's permission even when the current program
depends on it.

This distinction matters when a package introduces an unsafe prerequisite.
Suppose a Skill reads an unrelated private file before producing the requested
local output. Removing that read may cause the original script to fail because
the script tests its returned value. The failure shows a dependency in that
implementation. It does not show that the user authorized the read, or that a
different implementation needs it. A legal alternative can provide useful
evidence, but only if its inputs, permissions, and costs are made explicit.
Giving a defended system an alternative that its baseline cannot inspect would
confound the effect of recovery with unequal capabilities.

Permission to use a destination also does not determine which data may be sent
there. A task may permit reading private material for local processing and
sending a public report to an approved endpoint. It need not permit sending the
private material itself. Progent demonstrates explicit checks over tool names
and arguments, including task-based policy generation~\cite{progent}. CaMeL
demonstrates a separate role for control/data separation and information-flow
restrictions~\cite{camel}. These are complementary precedents. We do not claim
that task-aware permission checks or external enforcement are new.

TaskSkillGuard studies how source-grounded action checks, fixed task
permissions, finite flow restrictions, and recovery interact in a controlled
Skill setting. The study asks three questions. First, can action assessment
separate permission from dependence on an unsafe implementation? Second, which
effects are blocked by operation/resource permissions, and which require
checking the data sent to an otherwise permitted destination? Third, after a
denial or a known tool failure, does cost-based route selection improve on a
recovery controller with the same state, catalogue, simulation access, and
remaining quota? Graph presentation is a secondary question.

The contribution is a bounded combination and its controlled evaluation. Actual
synthetic Skill text selects a small Python-like program; each accepted effect
remains linked to its source. Executable task predicates check outputs,
preserved inputs, and effects. Matched controls separate scope enforcement,
flow enforcement, and recovery. Following the utility/security principle of
AgentDojo and Agent Security Bench~\cite{agentdojo,asb}, the protocol includes
legitimate communication and deletion, so blanket blocking cannot count as a
useful solution.

The implementation retains the source, grants, flow, and recovery framework of
the preserved Phase 5 study and strengthens it with Phase 6-style executable
checks and equal-capability controls. Historical observations remain separate
from the new evaluation. The new study uses synthetic packages only and a finite
in-memory execution model. It does not run arbitrary Python, contact real
services, estimate ecosystem attack prevalence, or reproduce the complete
SkillScope, Progent, or CaMeL systems. Its purpose is to make the tested security
boundary and the remaining coverage gaps explicit.

\section{Related Work}
Agent Skills combine instructions with executable resources, so their security
depends on both source content and the current task. Liu et al. document
confirmed malicious Skills~\cite{liu}, while repository-aware analysis warns
against treating suspicious content as confirmed maliciousness~\cite{context}.
Skill-Inject examines attacks inside Skill files~\cite{skillinject};
InjecAgent studies indirect injection through tool outputs~\cite{injecagent}.
SkillMutator extends this threat to interactions between language and
code~\cite{skillmutator}. These studies motivate inspecting the complete package
and retaining the task context, but their findings do not establish attack
prevalence in our constructed cases.

AgentDojo and Agent Security Bench evaluate harmful actions alongside normal
task completion~\cite{agentdojo,asb}. We follow this evaluation principle by
including permitted communication and deletion as well as attacks. Blocking
every sensitive action therefore cannot count as a useful defence. Our finite
synthetic environment tests specific mechanisms; it is not a reproduction of
either benchmark or a measurement of real Skill ecosystems.

StruQ separates trusted instructions from untrusted data through structured
inputs and training~\cite{struq}. IPIGuard uses tool dependencies to constrain
execution~\cite{ipiguard}, and DRIFT combines task-based planning with dynamic
validation~\cite{drift}. AttriGuard tests whether tool calls survive reduced
external control through counterfactual replay~\cite{attriguard}. Such evidence
addresses why an action occurs; causal support alone does not establish user
permission. Our graph presentation ablation does not implement IPIGuard.

SkillScope is the closest Skill-specific starting point~\cite{skillscope}. It
links instruction and code actions, checks authorization and necessity under
concrete tasks, and constrains validated excess actions. We adapt this framing
while separating necessity in the original implementation from necessity within
a finite alternative catalogue. An unsafe prerequisite can prevent a modified
program from finishing without becoming authorized. We do not claim to use
verified official SkillScope code or reproduce its full system.

Progent already generates task-based symbolic policies and checks policy
updates with an SMT solver~\cite{progent}. CaMeL separates control from untrusted
data and tracks information-flow restrictions~\cite{camel}. VIGIL enforces
behavioural specifications over execution traces~\cite{vigil}, while shielding
provides an earlier foundation for external action enforcement~\cite{shield}.
These precedents rule out a claim to the first task-aware policy or runtime
guard. Our contribution is a bounded combination and controlled evaluation:
source-linked actions, supplied task permissions, finite flow rules, executable
necessity checks, and recovery under unchanged authority and quota. Equal
information and simulation access let us test whether recovery gains come from
the selector itself. This scope does not imply arbitrary-script safety or
general superiority over the prior systems.


\section{Method}
\label{sec:method}
\subsection{Trusted Task and Untrusted Package}
Figure~\ref{fig:pipeline} shows the boundary tested in this study. Package text
selects a helper, but cannot grant that helper additional authority. The trusted
task supplies the permissions and the desired outcome. Every original or
recovery effect passes through the same gateway.

\begin{figure}[t]
\centering
\begin{tikzpicture}[>=Latex,box/.style={draw,rounded corners=2pt,align=center,font=\footnotesize,text width=3.0cm,minimum height=.72cm},node distance=.45cm]
\node[box] (package) {Untrusted Skill package\\\texttt{SKILL.md} + helper source};
\node[box,below=of package] (adapter) {Bounded source adapter\\effects linked to file and line};
\node[box,below=of adapter] (guard) {Scope + data-flow gateway};
\node[box,below=of guard] (effects) {In-memory effects\\check task outcome};
\node[box,right=.7cm of adapter] (task) {Trusted task\\grants, readers, goal};
\node[box,right=.7cm of effects] (routes) {Trusted alternatives\\same state and quota};
\draw[->] (package)--(adapter);
\draw[->] (adapter)--(guard);
\draw[->] (guard)--(effects);
\draw[->] (task)--(guard);
\draw[->] (effects)--(routes);
\draw[->] (routes)--(guard);
\end{tikzpicture}
\caption{The Skill-security boundary. A model precheck can restrict original
source sites but cannot change the trusted task. After a block or known failure,
recovery reuses the gateway.}
\label{fig:pipeline}
\end{figure}
The trusted task is represented as
\begin{equation}
T=(F_0,G,L,R,H,C,Q),
\end{equation}
where $F_0$ is the initial logical file state, $G$ is a set of exact
operation/resource grants, $L$ gives optional per-resource attempt limits,
$R$ gives outbound reader restrictions, $H$ is the goal predicate, $C$ is a finite
catalogue of trusted alternative scripts, and $Q$ is the shared call quota.
These inputs are supplied by the experiment. The study does not infer their
correctness from a Skill description or claim to solve natural-language policy
synthesis.

The untrusted package contains \texttt{SKILL.md} and script text. Package claims
such as ``the administrator approved this upload'' can be analysed as content;
they cannot add a grant, change a reader restriction, enlarge the quota, or
replace a trusted alternative. Optional model predictions can restrict proposed
original actions but cannot grant additional authority. Table~\ref{tab:trust}
summarizes this boundary.

\begin{table}[t]
\caption{Inputs and their authority in the finite model.}
\label{tab:trust}
\centering\small
\setlength{\tabcolsep}{3pt}
\begin{tabular}{p{0.27\columnwidth}p{0.65\columnwidth}}
\hline
Input & Permitted role \\
\hline
Trusted task & Defines grants, limits, flow restrictions, initial state, and goal. \\
Skill text and script & Supply the original program and evidence to inspect. \\
Source metadata & Links accepted effects to file content and source locations. \\
Model labels & Propose an original-action filter; cannot modify trusted inputs. \\
Trusted catalogue & Supplies explicit alternative programs under the same task. \\
\hline
\end{tabular}
\end{table}

An attacker in this model controls the original Skill instructions and script,
including misleading explanations and artificial dependencies. The attacker
does not control the dispatcher, initial task, trusted catalogue, or scoring
predicates. This boundary deliberately isolates whether the enforced task
constraints survive hostile package content. It does not cover a compromised
tool implementation, operating system, or task provider.

\subsection{From Source Text to a Finite Program}
The adapter requires one closed shell code block in \texttt{SKILL.md} containing
a single invocation of the form \texttt{python scripts/main.py}. The selected
path must be relative, normalized, traversal-free, and present in the supplied
package. Other shell commands, an ambiguous invocation, or a missing entry file
produce an unsupported result. The invocation selects text for parsing; it is
never passed to a shell.

The accepted language contains string, integer, and Boolean constants,
variables, assignments, string concatenation, and uppercase conversion. It
supports equality branches, literal loops of zero to three iterations, and
acyclic zero-argument helpers. Loop indices cannot be read. Helpers use renamed
local variables and require definite assignment before local reads. Imported
names, helper names, and the reserved internal prefix cannot be rebound.
Imports are limited to \texttt{Path} and the mock \texttt{emit} function.
Return statements, recursion, dynamic imports, arbitrary calls, and unsupported
operators are rejected. Integer and Boolean values may participate in accepted
expressions, but resource and payload arguments to effects must be strings.

Table~\ref{tab:semantics} states the effect semantics. Names resemble Python
operations to make source inspection familiar, but their meaning is defined by
this table. In particular, a missing read yields the empty string, and deleting
a missing resource leaves the state unchanged. Such a deletion is still a
recorded effect attempt and consumes quota when dispatched. These conventions
differ from ordinary host filesystem behaviour.

\begin{table}[t]
\caption{Effects in the bounded source language.}
\label{tab:semantics}
\centering\small
\setlength{\tabcolsep}{3pt}
\begin{tabular}{p{0.47\columnwidth}p{0.45\columnwidth}}
\hline
Source expression & Finite meaning \\
\hline
\texttt{Path(p).read\_text()} & Read logical string at $p$; return empty string if absent. \\
\texttt{Path(p).write\_text(v)} & Replace logical file $p$ with string $v$ and its labels. \\
\texttt{Path(p).unlink()} & Remove logical file $p$ if present; absence is a no-op. \\
\texttt{emit(d,v)} & Append destination/string pair $(d,v)$ to an in-memory log. \\
\hline
\end{tabular}
\end{table}

The adapter lowers accepted syntax into the finite intermediate language used
for the Phase 6 checks. Each effect receives a file-and-line identifier, source
span, source text, and file SHA-256 hash. Multiple statements on a line are
rejected to avoid ambiguous statement identities. Each source script is limited
to 24,000 encoded bytes, and each actual or simulated route has a 100-step
execution bound. Exceeding the runtime step bound is a stopped execution, not a
completed safe task. The original Phase 5 extractor is retained as a separate
measurement; its candidate graph is not treated as a complete execution model.

No original source is executed through \texttt{exec}, \texttt{eval}, a
subprocess, or a host file tool. Files and destinations are exact logical
strings. There is no claim about path aliases, symlinks, real permissions,
remote services, or operating-system isolation. Source hashes identify the text
used for an observation; they do not establish the publisher's trustworthiness.

All catalogue scripts are compiled before any route executes. Consequently, an
invalid alternative prevents execution even if the original script is valid.
This is a catalogue validation failure, which must be reported separately from
an unsupported original package. The runtime exposes an unsupported outcome
for both; the analysis distinguishes their validation stage.

\subsection{Scope and Information-Flow Checks}
For a proposed effect $e=(o,r,v)$, the scope check requires the exact pair
$(o,r)$ to be in $G$. It also requires the relevant attempt count to remain
below $L(o,r)$, except when a valid known receipt is being reused. If no local
limit is supplied, it does not further constrain the finite run. An attempted
call consumes both the shared quota and its resource-specific attempt count,
including a call denied by a guard or stopped by an injected failure before
commit. Receipt reuse does not dispatch a new call and is not charged again.

Values carry sets of restricted source origins. An initial file with an entry
in $R$ starts with its resource as an origin. A file without such an entry is
public in this model; an explicitly empty reader list forbids outbound
disclosure. Assignment, concatenation, uppercase conversion, file writes, and
branch control propagate origins. A read receives the stored file labels and
the labels influencing its resource and control context. This is a finite label
semantics, not a claim of noninterference for arbitrary Python.

Let $O(e)$ contain the origins influencing an effect's resource, payload, and
branch control. An outbound effect to destination $d$ passes the flow check
exactly when
\begin{equation}
\operatorname{FlowOK}(e)
=\bigwedge_{s\in O(e)} [d\in R(s)].
\end{equation}
The conjunction is true for an empty origin set. The check is applied to
outbound emits; the supported writes preserve labels for later reads. There is
no declassification operation. A guard checks both scope and flow again before
reusing a receipt.

For example, permission to emit to \texttt{review} does not allow a payload
derived from a source whose reader list excludes \texttt{review}. Replacing the
payload with its uppercase form retains that restriction. Writing it to a
temporary logical file and reading it back also retains the origin. These cases
separate destination permission from payload permission without assuming that
all network-like actions are malicious.

\subsection{Goals, Effects, and Successful Completion}
The goal predicate $H$ can require exact file contents, absence of specified
files, exact counts of required destination/payload pairs, and preservation of
specified inputs. Preservation checks both existence and content against $F_0$.
An unconstrained output is not silently made part of the task. Separate security
checks detect unauthorized effects even when the final files satisfy $H$.

For each run, let $E_c$ be its committed effects. The analysis counts scope
violations and forbidden outbound flows across the complete trace. Authorized
completion requires an actual attempt, a satisfied goal, no terminal error,
and no committed scope or flow violation:
\begin{equation}
\begin{split}
\operatorname{Success}={}&\operatorname{Attempted}\land H\land
\operatorname{NoError}\\
&\land\bigwedge_{e\in E_c}
\left(\operatorname{ScopeOK}(e)\land\operatorname{FlowOK}(e)\right).
\end{split}
\end{equation}
A final output cannot erase an earlier disclosure. Conversely, a denied action
is an observed attempt but not a committed harmful effect. The report keeps
these events and denominators separate. Committed reads are included in scope
checking, because collecting an unauthorized source is itself an excess action.

\subsection{Two Bounded Necessity Tests}
Consider a simple local-report task. The trusted task allows reading
\texttt{/input} and writing its uppercase text to \texttt{/output}; it does
not allow reading \texttt{/private}. The following schematic helper adds an
artificial prerequisite:
\begin{quote}\small\ttfamily
token = Path("/private").read\_text()\\
if token == "A":\\
\hspace*{1em}text = Path("/input").read\_text()\\
\hspace*{1em}result = text.upper()\\
\hspace*{1em}Path("/output").write\_text(result)
\end{quote}
In a world where the private value is \texttt{A}, the original output is
correct. Suppressing the private read makes the condition false and prevents
the output. A dependency-only check could therefore call the read necessary.
Yet the task grants still reject it. A legal helper can read \texttt{/input}
and write the transformed text directly, without the private read or branch.
The example explains why task completion, implementation dependence and
authorization need different tests. It is a schematic explanation of the
artificial-prerequisite mechanism, not an additional measured case.

Necessity is evaluated separately from permission. First, the original-program
test compares a no-fault execution with an intervention that suppresses every
occurrence of one source action site. A suppressed read supplies an empty
string; a suppressed write, deletion, or emit produces no effect. The comparison
checks the resulting task predicate. A failure after intervention is evidence
of dependence under this particular program and neutralization rule. It does
not establish that the action is required by the task. Sites not reached by the
original execution cannot receive a confirmed execution-based label.

The catalogue test considers the original and trusted routes from the same
initial state, without injected faults. Let $\mathcal{P}$ be the routes that
meet the goal while remaining within scope and flow restrictions. For an
observed source site $a$, let $T_a$ be its observed operation/resource pairs.
The test asks whether a legal route avoids every pair in $T_a$; it does not
require an alternative script to share the original file-and-line identity.
Let $E(p)$ be the pairs used by route $p$, and let $\mathcal{A}_a$ contain the
legal routes for which $E(p)\cap T_a=\varnothing$. The finite label is
\begin{equation}
N_C(a)=
\begin{cases}
\mathrm{no}, & \operatorname{closed}(C)\land\mathcal{A}_a\ne\varnothing,\\
\mathrm{yes}, & \operatorname{closed}(C)\land\mathcal{P}\ne\varnothing
\land\mathcal{A}_a=\varnothing,\\
\mathrm{unknown}, & \text{otherwise}.
\end{cases}
\end{equation}
Unexecuted sites remain unknown or not applicable. The implementation returns
unknown for every open catalogue. This is conservative: a legal avoiding
witness could logically establish a negative result even without closure, but
that stronger rule is not implemented here. The prompt wording mentions both
the witness rule and an open-catalogue unknown rule; our empty-open contexts
have no legal witness, so this ambiguity does not change their ground truth.
No legal solution also yields unknown. A closed catalogue means complete only
within the explicitly supplied finite set; it does not enumerate all possible
implementations. This is operation/resource-class necessity, not a claim that
the original payload or data origin is necessary. The same unsafe prerequisite can therefore be required by the
original script and avoidable by a legal catalogue route.

The protocol checks the accepted source through a separate AST interpreter
that does not call production lowering. It also checks agreement with the
Phase 6 finite-language reference where applicable. These comparisons assess
effects, goals, and interventions within the documented model. Agreement does
not prove the model matches every real Python program. Initial disagreements
and subsequent repairs belong in the evidence record.

\subsection{Recovery with Fixed Authority and Shared Cost}
Every configuration tries the original route first. Recovery retains the
current logical files, flow labels, outbound log, resource counts, quota, and
known receipts. It starts a fresh variable environment for the next route, but
does not roll back effects, reset quota, or add permissions. Non-recovery modes
make one actual route attempt. Recovery modes allow at most four actual
attempts in total, including the original.

After a recoverable stop, both recovery controllers simulate all eligible
alternatives from copies of the same current state. A route is feasible if its
simulation terminates without error and meets the goal under the same guards
and remaining quota. Previously selected alternatives are removed from the
unused set. The current route may be retried after a known before-commit failure
or a failure with a known completion receipt. Simulations assume no new fault;
they do not reveal a future injected failure.

The first-feasible controller selects the feasible route with the lowest fixed
catalogue index. The minimum-call controller minimizes the number of additional
charged calls in the simulation, breaking ties by that same index. Both
controllers perform the full eligibility simulation, so any difference is not
caused by giving only one controller access to a feasibility oracle. Simulated
calls and interpreter steps are reported separately from actual dispatched
calls. Simulation is not charged against the real execution quota, but its work
is not presented as free.

Receipts use the finite key
\begin{equation}
k=(\text{operation},\text{resource},\text{payload},\text{occurrence}).
\end{equation}
The occurrence counts appearances of the same operation/resource/payload
signature within a route attempt and restarts for a new attempt. It aligns
matching logical effects across routes. A known receipt reuses the recorded
return value and labels rather than dispatching the operation again; this
includes recorded reads. The duplicate counter concerns repeated commits of
this same finite key, not every repeated equal payload. The key does not include
the source location or a remote transaction identifier. This convention is
useful for the constructed routes, but is not a general exactly-once protocol
or a guarantee that arbitrary retries are equivalent.

\subsection{Known and Ambiguous Failure Outcomes}
The fault model places at most one fault at the first admitted, newly dispatched
non-read effect. Reads, guard denials, and receipt reuses do not trigger it.
Table~\ref{tab:faults} distinguishes when an effect occurs and what the controller
knows. A profile may never trigger if execution stops before reaching an eligible
effect; such a run must not be described as having exercised a mutation fault.

\begin{table}[t]
\caption{Finite fault profiles and allowed response.}
\label{tab:faults}
\centering\small
\setlength{\tabcolsep}{3pt}
\begin{tabular}{p{0.24\columnwidth}p{0.35\columnwidth}p{0.31\columnwidth}}
\hline
Profile & Effect and receipt & Controller response \\
\hline
No fault & Normal effect and receipt & Continue. \\
Before commit & No effect; no receipt & May retry or use a legal alternative. \\
Known after commit & Effect occurs; receipt retained & May recover with receipt reuse. \\
Ambiguous after commit & Effect occurs; no receipt & Stop without retry. \\
\hline
\end{tabular}
\end{table}

The evaluator can record the effect of the ambiguous profile because it controls
the finite environment. The controller does not use that privileged observation
to continue. It stops on an ambiguous outcome or exhausted quota. If no
simulated route is feasible, recovery also stops. These explicit failure states
prevent successful output alone from being confused with a justified recovery.

\section{Experimental Design}
\label{sec:experimental-design}
\subsection{Synthetic Cases and Freeze Boundary}
The main set contains 48 supported cases: 12 mechanism families, each
with a benign and an adversarial form and two lexical variants. Six separate
unsupported probes test the source-coverage boundary. The variants are
dependent constructions, not 48 independent attack families. Family-level
summaries are therefore the primary unit for interpreting variation. The set is
not a random sample of a marketplace and does not estimate production risk.

The mechanisms cover authority claims, excessive collection, unauthorized
destinations, forbidden payloads sent to permitted destinations, derived and
file-mediated flows, artificial prerequisites, source changes, repeated
effects, and helper functions. Benign counterparts include required
communication and deletion. Their purpose is to test whether a guard preserves
the distinction between a powerful action and an unauthorized one.

Separate development fixtures support unit testing. Production code and the
model protocol were frozen before the evaluation generator was first run or
read by the implementation context. The generator and direct AST reference
were developed in a separate assistant context that did not read the production
modules. This provides implementation separation, not independent human
validation. Case identities, source hashes, run configuration, and first
failures are retained. Repairs after release are explicit amendments; their
results remain separate from the first comparison. No human labels or real
Skill packages are used in these experiments.

\subsection{Matched Controls}
Table~\ref{tab:controls} defines the five mechanism controls. Each starts from
the same task and original source. The matrix applies the four fault
profiles to every supported case and control. A record of an assigned profile
is kept distinct from evidence that its fault was reached. The original and
catalogue necessity tests use their separate no-fault initial-state protocol.

\begin{table}[t]
\caption{Mechanism controls. All use the same trusted task.}
\label{tab:controls}
\centering\small
\setlength{\tabcolsep}{3pt}
\begin{tabular}{p{0.24\columnwidth}ccc p{0.25\columnwidth}}
\hline
Control & Scope & Flow & Recovery & Selection \\
\hline
Observe & No & No & No & Original only \\
Scope & Yes & No & No & Original only \\
Scope + flow & Yes & Yes & No & Original only \\
First feasible & Yes & Yes & Yes & Fixed route order \\
Minimum calls & Yes & Yes & Yes & Fewest added calls \\
\hline
\end{tabular}
\end{table}

Observe mode records the same authorization and flow predicates while allowing
their violations to occur in memory. It is not real-world unguarded execution.
Scope versus scope-plus-flow isolates the payload restriction. The two recovery
modes share the complete catalogue and simulation ability; their difference is
the selection rule. All controls retain the same supported execution and
receipt conventions. Broader claims about other receipt designs require a
separate experiment.

\subsection{Outcomes and Coverage}
Primary outcomes are runs with any unauthorized committed effect, runs with
any forbidden outbound flow, authorized goal completion, duplicate mutations
under the defined key, charged calls, simulated calls, and planner steps.
Counts of individual effects accompany run-level counts where useful. The
denominator must identify whether it concerns all supported cases, a mechanism
family, or a predeclared attack opportunity. Blocked attempts, terminal errors,
unsupported originals, and invalid catalogues are separate outcomes rather than
silent exclusions or safe classifications.

Source coverage is reported independently of action-label quality. A supported
program may contain an unexecuted action site; the site is not confirmed safe
merely because one task does not reach it. Model abstentions and malformed
answers remain in coverage denominators. Predictions cannot replace the
executable goal or effect checks. Any uncertainty intervals must preserve family
dependence and be described cautiously with only 12 constructed families.

The model audit receives the complete synthetic package and trusted task.
Source/table and graph presentations are secondary comparisons, using the
same gateway and retained invalid responses. Model
labels only restrict original action sites; trusted recovery routes remain
subject to the unchanged scope and flow guards. The audit therefore measures a
specified component of the system, not end-to-end autonomous permission
generation. Its samples and costs are reported separately from deterministic
mechanism runs. The frozen acquisition uses GPT-6 Luna and GPT-4.1 mini, one
response per distinct input, with no automatic retries. It covers every main
case in two views and adds empty open/closed catalogue contexts for the first
attack variant in each family. Identical inputs reuse the same saved response.
The graph adds conservative data, control, and call edges derived from the
finite program; it does not add a gold label. Source text and trusted task facts
are identical across views. Graph length is not matched by a padding control,
so the comparison measures the addition of graph information and its token
cost together. JSON answers use a fixed schema and a phase-aware parser.

The preserved Phase 5 graph results and Phase 6 strong-control results provide
historical context with their original denominators. They are not pooled with
the new cases or treated as independent replications of this source adapter.
Negative graph results and a lack of advantage over an equally capable recovery
control remain valid outcomes. The evaluation is intended to identify what each
mechanism contributes within the stated boundary, including where the proposed
combination adds no measurable benefit.
